\section{Conclusion}
\label{sec:conclusion}

We presented the first systematic decomposition of the R-matrix
mechanism in Hierarchical Multi-Label Classification. Our 4-way ablation
across three hierarchy types reveals that \textbf{inference-time
bottom-up reconciliation accounts for over 90\% of the R-matrix
benefit}, while training-time consistency loss provides negligible
additional value. This finding simplifies HMC system design:
practitioners should always apply reconciliation at inference (zero
cost, maximum benefit) and can safely omit the training-time MC-loss.

Our scaling law analysis---the first of its kind for HMC---shows that
the R-matrix benefit follows $\Delta\text{F1} \propto D \cdot
(N/C)^{-0.5}$: deeper hierarchies and data-scarce regimes benefit most.
This provides a practical heuristic for when to invest in hierarchical
constraints.

Together with the HMC-Torch platform~\cite{sette2026hmctorch}, these
findings provide a foundation for principled HMC system design:
reconciliation always on, consistency loss never, and hierarchical
constraints deployed where the scaling law predicts they will help most.

\subsection*{Reproducibility}

All experiments use fixed seeds and are fully reproducible via the
HMC-Torch platform (\url{https://github.com/sette/hmc-torch}). The
ablation and scaling law experiment scripts are included in the
repository under \texttt{experiments/run\_ablation.py} and
\texttt{experiments/run\_scaling\_laws.py}. Multi-seed results with 5
seeds and standard deviations are provided for all primary benchmarks.
